Jedes der $N$ Heliumatome hat die mittlere kinetische Energie $\frac{3}{2}kT$. Die gesamte kinetische Energie ist $E = N\cdot\frac{3}{2}kT = \frac{3}{2}nRT$ (mit $N = nN_\text{A}$ und $R = N_\text{A}k$). Steigt die Temperatur um $\Delta T$, nimmt sie um $\Delta E = \frac{3}{2}nR\,\Delta T$ zu.

Die Stoffmenge ist mit $M = \MHeO$ (FoTaBe S.~228)
\begin{align}
 n &= \nHeF = \frac{\mHeO}{\MHeO} = \nHe
\end{align}
und die zuzuführende Energie
\begin{align}
 \Delta E &= \dEF \\
   &= \frac{3}{2}\cdot\nHe\cdot\ncR\cdot\dTh \\
   &= \dE \approx \result{\dEP{3}{2}}
\end{align}
